% BITS-Bench: ICML submission draft.
% To build with the official template, replace the \documentclass/\usepackage block below with the
% ICML style files (icml20XX.sty, icml20XX.bst) and \usepackage[accepted]{icml20XX} as instructed in
% the call for papers. Everything marked \todo{} needs LLM results that `scripts/run_grid.sh` produces.
\documentclass[11pt]{article}
\usepackage[margin=1in]{geometry}
\usepackage{amsmath,amssymb,amsthm}
\usepackage{booktabs}
\usepackage{graphicx}
\usepackage{hyperref}
\usepackage{xcolor}
\usepackage{natbib}
\usepackage{enumitem}

\newtheorem{proposition}{Proposition}
\newtheorem{definition}{Definition}
\newcommand{\todo}[1]{\textcolor{red}{[TODO: #1]}}
\newcommand{\Hs}{\mathcal{H}}
\newcommand{\X}{\mathcal{X}}
\newcommand{\Y}{\mathcal{Y}}
\newcommand{\E}{\mathbb{E}}
\newcommand{\bench}{\textsc{BITS-Bench}}

\title{\bench: Separating Information Gathering from Inference\\in Language-Model Discovery Agents}
\author{Anonymous Authors}
\date{}

\begin{document}
\maketitle

\begin{abstract}
Language-model agents are increasingly asked to do science: choose experiments, interpret their
outcomes, and arrive at a correct theory. When such an agent fails, it is usually impossible to say
whether it ran the wrong experiments or drew the wrong conclusions from the right ones, because
end-to-end benchmarks have no ground truth for what an ideal reasoner could have concluded from the
agent's own data. We introduce \bench{}, a benchmark of active rule discovery in which every hidden
rule is drawn uniformly from a finite, fully enumerated hypothesis space of up to $3\times10^5$
semantically distinct rules across four domains (number concepts, Boolean functions, regular
languages and list programs), with budgets of 8--24 experiments and an optional noisy-oracle regime.
Exact Bayesian inference over these spaces yields, for every agent and every checkpoint, the
information its experiments gathered in bits and the accuracy of the Bayes-optimal learner given
those same experiments. The agent's shortfall from a greedy expected-information-gain reference then
decomposes \emph{exactly} into a \emph{gathering regret} and an \emph{inference regret}, the latter
non-negative in expectation, both for prediction and for exact identification of the rule. With
perfect inference, random experimentation recovers most of the reference's predictive accuracy but
identifies the rule only 15--69\% of the time on binary domains, versus 87--100\% for the reference:
the identification track is where experimental design matters. \todo{One or two sentences of headline
LLM findings: e.g.\ fraction of total regret due to inference vs gathering for frontier models; rate of
logically entailed (zero-information) experiments; whether choosing their own experiments helps
models relative to random data.}
\end{abstract}

\section{Introduction}

A recurring hope for capable AI systems is that they will accelerate research: propose hypotheses,
design experiments, and revise beliefs in light of results. Recent benchmarks evaluate this end to
end. ASI-Bench \citep{asibench2026} grades agents on project-level research tasks under decreasing
methodological guidance and finds that performance drops sharply when agents must choose methods
themselves; DiscoveryWorld \citep{jansen2024discoveryworld}, BoxingGym \citep{gandhi2025boxinggym},
ACED-Bench \citep{acedbench2026} and ActiveSciBench \citep{autoscilab2026} place agents in simulated
environments where they must experiment to discover hidden laws. These benchmarks establish that
current agents are weak experimenters. They do not establish \emph{why}.

Two failure modes are confounded in any end-to-end score. An agent may \emph{gather} too little
information (redundant, confirmatory, or uninformative experiments), or it may \emph{use} the
information it gathered poorly (failing to rule out hypotheses its data contradicts, or over-committing
to ones it does not support). The remedies differ (better exploration policies versus better
inductive inference), so the distinction matters for both evaluation and method development. Making it
requires knowing what an ideal reasoner could have concluded from the agent's \emph{own} data. In
open-ended environments this is unknowable.

\bench{} makes it knowable by construction. Each task is a black box computing a rule drawn uniformly at
random from a finite hypothesis class that we enumerate completely and tell the agent about precisely.
The agent chooses inputs, observes outputs, and periodically predicts the rule's outputs on held-out
inputs; at the end it states the rule. Because the hypothesis space is explicit, the exact posterior
after any experiment sequence is computable, and with it the Bayes-optimal predictions, the information
gained in bits, and the expected information gain of every experiment the agent \emph{could} have run.

\paragraph{Contributions.}
\begin{enumerate}[leftmargin=*,itemsep=2pt]
\item A benchmark of ten configurations over four formal domains with hypothesis spaces of $2^{4.7}$ to
$2^{18.3}$ rules, a noisy-oracle variant, procedurally generated episodes with a private test salt, and
three conditions (active, passive-random, passive-EIG) (\S\ref{sec:bench}).
\item An exact decomposition of an agent's regret relative to a strong reference experimenter into
gathering and inference components, for both prediction and identification, with inference regret
provably non-negative in expectation (\S\ref{sec:metrics}).
\item Per-experiment diagnostics that are only available with exact posteriors: the information gain of
each chosen experiment relative to the best available one, and the rate of \emph{entailed}
experiments whose outcome was already determined by earlier results.
\item Reference results showing that predictive accuracy is a poor discriminator of experimental
design while exact identification is a sharp one (\S\ref{sec:refs}), and an evaluation of
\todo{N} language models \todo{(\S\ref{sec:llm})}.
\end{enumerate}

\section{Related work}

\paragraph{Optimal experimental design and active learning.} Choosing experiments to maximise expected
information gain dates to \citet{lindley1956}; see \citet{chaloner1995} and \citet{rainforth2024} for
reviews. For noise-free identification over a finite hypothesis class, the greedy most-informative
query is generalised binary search, whose expected number of queries is within a logarithmic factor of
optimal \citep{dasgupta2004greedy,golovin2011adaptive}. Version spaces \citep{mitchell1982} are the
noise-free posterior supports we track.

\paragraph{Human hypothesis testing.} Our number domain extends Tenenbaum's number game
\citep{tenenbaum1999bayesian} with an explicit hypothesis class. People choosing tests famously prefer
cases they expect to confirm their current hypothesis \citep{wason1960,klayman1987confirmation};
information-gain models of human question asking are reviewed by \citet{coenen2019asking}. Our entailed-experiment
rate measures an extreme form of this bias: tests whose outcome is logically determined.

\paragraph{LLM information gathering.} Twenty-questions-style evaluations and methods
\citep{bedllm2025,asig2026} study question asking with EIG rewards; WILT \citep{banatt2024wilt} adapts
the Wason 2-4-6 task; QuestBench \citep{li2025questbench} evaluates asking for missing information.
BoxingGym \citep{gandhi2025boxinggym} also uses EIG, estimated by sampling, to score experimental design
in probabilistic-model environments. \bench{} differs in computing posteriors \emph{exactly} over
enumerated spaces, which is what makes the regret decomposition exact and the inference term provably
non-negative, and in scoring identification and prediction on the same episodes.

\paragraph{Scientific-discovery agents.} ASI-Bench \citep{asibench2026}, DiscoveryWorld
\citep{jansen2024discoveryworld}, DiscoverPhysics \citep{discoverphysics2026}, ACED-Bench
\citep{acedbench2026} and HypoArena \citep{hypoarena2026} evaluate open-ended discovery. These are
complementary: they are realistic but unable to attribute failure; \bench{} is formal but fully
attributable. Program-induction benchmarks such as list functions \citep{rule2020child} and ARC
\citep{chollet2019measure} evaluate passive induction from given examples; our passive conditions are
of this kind.

\section{The benchmark}\label{sec:bench}

\paragraph{Setup.} A configuration fixes an input space $\X$, an output space $\Y$, a finite hypothesis
class $\Hs \subseteq \Y^{\X}$ of semantically distinct functions, an experiment budget $B$, and an
observation-noise rate $\epsilon$. An episode draws a target $h^\star \sim \mathrm{Unif}(\Hs)$ and a
set of $n$ test inputs from a fixed distribution $D$ over $\X$, independently of $h^\star$. The agent is
given an exact natural-language specification of $\Hs$ and of the uniform prior. At steps
$t=1,\dots,B$ it chooses $x_t\in\X$ and observes $y_t = h^\star(x_t)$, flipped with probability
$\epsilon$ in noisy binary configurations. After $k\in\{B/8,B/4,B/2,B\}$ experiments it predicts
$h^\star$ on the test inputs; after $B$ it also states a rule.

\paragraph{Domains.} Table~\ref{tab:refs} lists the ten configurations. \emph{Numbers}: $\X=\{1..100\}$,
rules are one atomic property (residues modulo 2--10, thresholds, primality, squares, digit
membership, digit-sum parity and thresholds, tens digit; each possibly negated) or the AND/OR of two.
\emph{Boolean}: $\X=\{0,1\}^8$; conjunctions/disjunctions of $\le 3$ literals, extended with 2-term DNF
and 2-clause CNF. \emph{Regular}: strings over $\{a,b\}$ of length 1--8; languages of DFAs with $\le 3$
or $\le 4$ states. \emph{Lists}: lists of 1--6 digits; pipelines of up to three of 26 primitives
(sorting, filtering, arithmetic modulo 10, reductions, \dots). Binary rules are restricted to positive
rates in $[0.1,0.9]$. Hypotheses that agree on all of $\X$ (for lists, on all lists of length $\le 3$
plus 400 fixed longer lists) are merged, keeping the shortest description, so the prior is uniform
over distinct functions. Sizes are pinned by regression tests.

\paragraph{Protocol details.} Each checkpoint is answered in a fork of the conversation (the history
plus one turn) that is then discarded, so the test inputs never enter the history from which the agent
experiments and cannot be queried. Every message restates all results so far, so the benchmark does not
test long-context recall. Malformed experiments may be retried twice before the step is forfeited.
Submitted rules are compiled (restricted Python expressions, regular expressions or DFAs, DSL
pipelines) and counted correct iff equivalent to $h^\star$ on all of $\X$.

\paragraph{Conditions.} In \emph{active} episodes the agent chooses its experiments. In
\emph{passive-random} and \emph{passive-EIG} episodes it sees the first $k$ results of the random or
greedy-EIG reference experimenter and only predicts. Since experiments are chosen from past
observations only, the posterior is the same however they were chosen, so all three conditions share
the same Bayes-optimal reference for a given dataset.

\paragraph{Splits.} Episodes are generated deterministically from a hash of
(split, salt, domain, level, $\epsilon$, index). The dev split (30 episodes per configuration; 390 with
the noisy suite) uses a public salt. The test split uses a private salt, so fresh instances can be
drawn for every evaluation and memorising released episodes does not help.

\section{Measuring gathering and inference}\label{sec:metrics}

Let $D_k = ((x_1,y_1),\dots,(x_k,y_k))$ be the agent's first $k$ observations and
$p(h\mid D_k)\propto \prod_{t\le k} p(y_t\mid h(x_t))$ the exact posterior under the uniform prior.

\begin{definition}[Information and Bayes accuracy]
The information gathered is $I(D_k) = \log_2|\Hs| - H[p(\cdot\mid D_k)]$ bits. The Bayes predictor
outputs $\hat y^{B}_i = \arg\max_y p(h(x_i)=y\mid D_k)$ on test input $x_i$; its realised accuracy is
$\mathrm{acc}^B(D_k)$.
\end{definition}

Let $R_k$ denote the accuracy of the Bayes predictor on the data collected by a reference experimenter
in $k$ steps (averaged over reference runs). For an agent with realised accuracy $\mathrm{acc}_k$:
\begin{equation}\label{eq:decomp}
\underbrace{R_k - \mathrm{acc}_k}_{\text{total regret}}
= \underbrace{R_k - \mathrm{acc}^B(D_k)}_{\text{gathering regret}}
+ \underbrace{\mathrm{acc}^B(D_k) - \mathrm{acc}_k}_{\text{inference regret}} .
\end{equation}

\begin{proposition}\label{prop:infer}
If the agent's predictions depend on $h^\star$ only through $D_k$, then
$\E[\mathrm{acc}^B(D_k) - \mathrm{acc}_k \mid D_k, x_{1:n}] \ge 0$.
\end{proposition}
\begin{proof}
Given $D_k$ and the test inputs (drawn independently of $h^\star$), $h^\star \sim p(\cdot\mid D_k)$
because the target is drawn from the stated prior. Any prediction $\hat y_i$ that is conditionally
independent of $h^\star$ is correct with probability $p(h^\star(x_i)=\hat y_i\mid D_k) \le
\max_y p(h^\star(x_i)=y\mid D_k)$, which is the probability that $\hat y^B_i$ is correct. Average over
$i$.
\end{proof}
The forked checkpoints guarantee the condition: the agent never sees test labels. Gathering regret, in
contrast, can be negative, since greedy EIG is not optimal in every episode.

\paragraph{Identification.} The Bayes-optimal guess of the rule is the MAP hypothesis, correct with
probability $m(D_k) = \max_h p(h\mid D_k)$. With $M_B$ the reference's MAP mass and $c\in\{0,1\}$
whether the agent's stated rule is exactly right,
$M_B - c = (M_B - m(D_B)) + (m(D_B) - c)$, and $\E[m(D_B) - c\mid D_B]\ge 0$ by the same argument.

\paragraph{Scores.} \emph{Predict} is accuracy above the prior baseline as a fraction of the
reference's, $(\overline{\mathrm{acc}} - \overline{\mathrm{acc}}_0)/(\overline{R} -
\overline{\mathrm{acc}}_0)$, averaged over checkpoints (an area under the learning curve).
\emph{Discover} is $\bar c / \bar M_B$. Both are computed per configuration and averaged with equal
weight; the headline \bench{} score is their mean. We also report \emph{information efficiency}
$\sum I(D_k)/\sum I(D^{\mathrm{ref}}_k)$ and \emph{design efficiency}
$(I - I^{\mathrm{rand}})/(I^{\mathrm{ref}} - I^{\mathrm{rand}})$, negative when an agent's experiments are
less informative than random ones.

\paragraph{Per-experiment diagnostics.} For each experiment $x_t$ we compute its expected information
gain $\mathrm{EIG}_t(x_t)$ under $p(\cdot\mid D_{t-1})$ and the best available
$\max_x \mathrm{EIG}_t(x)$. Their ratio measures myopic design quality. In noise-free configurations an
experiment with $\mathrm{EIG}_t(x_t)=0$ while the rule is still undetermined is \emph{entailed}: its
outcome was already implied by earlier results. On the noise-free dev configurations, random
experimenters spend 16\% of their open-question experiments on entailed tests and attain 0.66 of the
best available information gain per experiment; greedy EIG runs none.

\section{Reference results}\label{sec:refs}

\begin{table}[t]
\centering
\small
\begin{tabular}{lrrrrrrrrr}
\toprule
config & $\log_2|\mathcal{H}|$ & $B$ & prior & \multicolumn{2}{c}{accuracy (AUC)} & \multicolumn{2}{c}{bits at $B$} & \multicolumn{2}{c}{identified at $B$} \\
& & & & random & EIG & random & EIG & random & EIG \\
\midrule
numbers/L1 & 7.6 & 16 & 0.50 & 0.86 & 0.89 & 6.8 & 7.6 & 0.69 & 1.00 \\
numbers/L2 & 14.1 & 16 & 0.45 & 0.79 & 0.82 & 9.8 & 13.7 & 0.15 & 0.87 \\
boolean/L1 & 10.1 & 16 & 0.46 & 0.86 & 0.89 & 7.9 & 10.1 & 0.38 & 0.99 \\
boolean/L2 & 13.5 & 16 & 0.44 & 0.73 & 0.77 & 11.1 & 13.5 & 0.30 & 1.00 \\
boolean/L3 & 18.3 & 24 & 0.44 & 0.78 & 0.82 & 15.0 & 18.2 & 0.26 & 0.99 \\
regular/L1 & 9.8 & 16 & 0.55 & 0.80 & 0.83 & 8.4 & 9.8 & 0.56 & 1.00 \\
regular/L2 & 15.6 & 16 & 0.55 & 0.75 & 0.78 & 11.7 & 15.5 & 0.24 & 0.98 \\
lists/L1 & 4.7 & 8 & 0.14 & 0.94 & 1.00 & 4.7 & 4.7 & 1.00 & 1.00 \\
lists/L2 & 8.8 & 8 & 0.10 & 0.83 & 0.94 & 8.7 & 8.8 & 0.97 & 1.00 \\
lists/L3 & 12.5 & 8 & 0.16 & 0.72 & 0.88 & 12.0 & 12.5 & 0.85 & 1.00 \\
numbers/L1 $\epsilon$=0.1 & 7.6 & 16 & 0.39 & 0.80 & 0.83 & 5.0 & 6.7 & 0.45 & 0.82 \\
boolean/L1 $\epsilon$=0.1 & 10.1 & 16 & 0.51 & 0.81 & 0.83 & 5.0 & 8.0 & 0.17 & 0.63 \\
regular/L1 $\epsilon$=0.1 & 9.8 & 16 & 0.54 & 0.71 & 0.74 & 6.3 & 7.8 & 0.36 & 0.65 \\
\bottomrule
\end{tabular}

\caption{Reference experimenters with exact Bayesian inference on the dev split (30 episodes per
configuration). Accuracy is the area under the learning curve over the four checkpoints; ``identified''
is the posterior mass of the MAP rule after the full budget $B$.}
\label{tab:refs}
\end{table}

Table~\ref{tab:refs} shows two regimes. On the binary domains, random experiments with perfect
inference achieve predictive accuracy within 0.03--0.04 of greedy EIG, because a random input to a
balanced rule already yields close to one bit. They gather substantially fewer bits by the end of the
budget, however (e.g.\ 9.8 vs 13.7 of 14.1 on numbers/L2), and the last few bits are exactly what
separates a nearly-right theory from the right one: random experimentation identifies the rule in
15--69\% of episodes versus 87--100\% for the reference. On list programs, where each output carries many
bits, both experimenters identify the rule almost always and the remaining difficulty is inferential.
Observation noise lowers identification for both and widens the gap in bits. This motivates reporting
prediction and identification separately: a benchmark that scored only prediction would barely
distinguish a random experimenter from a near-optimal one.

\section{Evaluating language models}\label{sec:llm}

\paragraph{Models and settings.} \todo{List models (frontier models from several providers at two
reasoning-effort levels, smaller models from the same families, and open-weight models via vLLM), reasoning
settings, max output tokens, number of episodes; all runs use \texttt{scripts/run\_grid.sh}.}

\paragraph{Main results.} \todo{Leaderboard table from \texttt{reports/dev/leaderboard.tex}: BITS score,
Predict, Discover, gathering/inference regrets for both tracks, information and design efficiency,
entailed-experiment rate. Learning-curve figure \texttt{curves.png} and decomposition figure
\texttt{decomposition.png}.}

\paragraph{Analyses.}
\begin{itemize}[leftmargin=*,itemsep=2pt]
\item \textbf{Where does the regret come from?} \todo{Share of total regret that is inference vs.\
gathering, per domain and per model.}
\item \textbf{Does choosing experiments help the model?} \todo{Compare active vs passive-random accuracy
for each model. For the Bayes learner, active-EIG beats random by the margins in Table~\ref{tab:refs};
does the same hold for models, or do they learn more from random data than from their own?}
\item \textbf{Can models exploit ideal data?} \todo{Passive-EIG inference regret: how much of the
information in near-optimal experiments the model uses.}
\item \textbf{Scaling with hypothesis entropy.} \todo{Regret vs.\ $\log_2|\Hs|$ across configurations.}
\item \textbf{Noise.} \todo{Noisy vs noise-free L1; do models repeat experiments to average out noise?}
\item \textbf{Confirmation and entailed tests.} \todo{Entailed-experiment rate and EIG ratio per model;
qualitative examples from transcripts.}
\item \textbf{Reasoning effort.} \todo{Effort sweep: which regret component additional thinking
reduces.}
\item \textbf{Human baseline.} \todo{Optional: \texttt{bitsbench play} on a subset of episodes.}
\end{itemize}

\section{Limitations}

The rule families are formal and small by design, since exact posteriors require enumerable spaces.
\bench{} isolates a precondition for automated research (efficient, correct learning from self-chosen
experiments) and does not measure research ability as such. Greedy EIG is a strong reference but not an
upper bound on information gathered; the decomposition's guarantee concerns the inference term. Telling
the agent the exact hypothesis class and prior is a strong simplification of real science, where
hypothesis-space construction is itself the hard part \citep{hypoarena2026}; we view \bench{} as
complementary to benchmarks that test it.

\bibliographystyle{plainnat}
\bibliography{refs}

\appendix
\section{Prompts}
\todo{Include \texttt{bitsbench prompt <domain> <level>} output for each configuration and the
checkpoint prompt.}

\section{Hypothesis-space construction}
Spaces are built by enumerating each grammar in order of description length, evaluating every
hypothesis on the full input space (for lists, the reference pool), discarding unbalanced binary rules,
and keeping the first hypothesis of each equivalence class. For lists, prefixes that agree on the pool
are pruned before extension, which is exact because extensions of equal prefixes are equal. For regular
languages, all DFAs with $k\le 4$ states and start state $q_0$ are simulated on all 510 strings with a
prefix trie.

\end{document}
